\formulary{Completorium commune}
\label{form:commune}

Converte\footcite[14r]{diu1523} nos Deus salutaris noster.
Et averte iram tuam a nobis.\\
Deus in adiutorium \editorial{meum intende}.
Domine ad adiuvandum me festina.\\
Gloria Patri \editorial{et Filio et Spiritui Sancto.
  Sicut erat in principio et nunc et semper et in saecula saeculorum amen.}
Alleluia.

\rubrica{Psalmus}\label{psalmi}
\input{psalmi/ps4.tex}
Gloria Patri et.

\rubrica{Psalmus}
\input{psalmi/ps30i.tex}
Gloria Patri et Filio.

\rubrica{Psalmus}
\input{psalmi/ps90.tex}
Gloria Patri et Filio.

\rubrica{Psalmus}
\input{psalmi/ps133.tex}
Gloria Patri et.

\rubrica{An\editorial{tiphona}.}
Alleluia.\footnote{%
  \rubrica{\cite[121r]{bp1502}: Ad completorium psalmi consueti quos sequitur \incipit{Alleluia.}}}
\edissue{Not sure what exactly the alleluia means. An antiphon super psalmos? But XIV A 19 doesn't notate it. diu1523 rubricates it as antiphon. I vaguely remember that historians of the office mention an ancient practice of saying psalms with alleluia, predating the antiphonal psalmody. Read more on this topic.}

\rubrica{Capitulum}
Pacem et veritatem
atque iustitiam\footnote{%
\cite[121r]{bp1502}: \enquote{atque iustitiam} deest
}
diligite: ait Dominus omnipotens.
Deo \editorial{gratias}.

\rubrica{Hymnus}
\edissue{bp1502 has in the formulary Salvator mundi
  and Te lucis mentions only in a rubric afterwards -- explain why}
\hymnus{te_lucis}

\rubrica{V.}\label{commune:versus:custodi}
Custodi nos Domine ut pupillam oculi.
Sub umbra allarum tuarum protege nos.\quickref{can:commune_versiculus}

\rubrica{Ad Nunc dimittis antiphona}
Salva nos Domine vigilantes:
custodi nos dormientes:
ut vigilemus in Christo
et requiescamus in pace.\quickref{can:commune_salva1}

\rubrica{\editorial{Canticum Symeonis}}\footnote{%
  Canticum in \cite{diu1523} deest,
  desumptum est psalterio \cite{ps1502}}
\input{psalmi/nuncdimittis}

\rubrica{\editorial{Preces}}
\label{preces}
\chapter{Preces}
\label{preces}

Kyrieleison. Kyrieleison. Kyrieleison.\\
Christeleison. Christeleison. Christeleison.\\
Kyrieleison. Kyrieleison. Kyrieleison.\edissue{compare other contemporary local sources how kyrie was most usually written}
\\
Pater noster. Et ne nos.\edissue{are these incipits of recited portions, or is the rest said silently?}\\
In pace in idipsum. Dormiam et requiescam.\\
Credo in Deum. Carnis resurrectionem. Et vitam.
\edissue{are these incipits of recited portions, or is the rest said silently? Clarify the execution, provide full texts.}
\\
Benedicamus Patrem et Filium cum Sancto Spiritu.
Laudemus et superexaltemus sum in saecula.\\
Benedictus es, Domine, in firmamento coeli.
Et laudabilis et gloriosus et superexaltatus in saecula.\\
Benedicat nos omnipotens Dominus. Amen.\edissue{uncertain}
\\
Dignare Domine nocte ista.
Sine peccato nos custodire.\\
Miserere nostri, Domine, miserere nostri.\\
Fiat misericordia tua, Domine, super nos.
Quemadmodum speravimus in te.\\
Sacerdotes tui induantur iustitia.\edissue{verify exact spelling, this one differs from now standard Vulgate}
Et sancti tui exsultent.

\noindent\rubrica{\editorial{Versus proprius tempori aut festo hic adnectitur.}}

\noindent{}Domine exaudi orationem meam.
Et clamor meus \editorial{ad te veniat}.\\
Dominus vobiscum.\edissue{This dialogue is probably a standard introduction for the prayer, not part of the preces as such.}
\editorial{Et cum spiritu tuo.}

\footcite[1v]{xiv_a_19}
\edissue{Try to find in printed breviaries -- so far no luck}


\rubrica{Oremus.}
Vigila super nos aeterne Salvator:
ne nos apprehendat ille callidus tentator:
quia tu nobis factus es: sempiternus adiutor.

\rubrica{\editorial{Alia}}\label{commune:oratio:visita}
Visita quaesumus Domine habitationem istam:
et omnes insidias inimici ab ea longe repelle:
et angeli tui sancti\footnote{\cite[121r]{bp1502}: sancti tui}
in ea habitantes nos in pace custodiant:
et benedictio tua sit super nos semper.\footnote{\cite[121r]{bp1502}: sit semper nobiscum.}

\rubrica{Alia oratio}\label{commune:oratio:salvator}
Salvator mundi salva nos in hac nocte et in omni tempore:
et intercedente pro nobis
beata et gloriosa semperque Virgine Dei Genitrice Maria\footnote{%
  \cite[121r]{bp1502}: beata Dei Genitrice et Virgine Maria}
cum omnibus sanctis tuis:
pacem tuam nostris concede temporibus:
et omnia mala a nobis propiciatus
averte.\footnote{\cite[121r]{bp1502}: exclude}
Qui cum Deo Patre et Spiritu Sancto vivis et regnas Deus.
Per omnia saecula saeculorum.

\rubrica{\editorial{Conclusio}}
Dominus vobiscum.
Et cum spiritu tuo.\\
Benedicamus Domino. Deo gratias.\edissue{notation}
\footnote{%
  \cite[121r]{bp1502}:
  \rubrica{Supra notatum capitulum et hymnus
  \incipit{Te lucis,}
  versus et antiphonae et orationes dicuntur ad completorium
  per circulum anni,
  praeter Quadragesimam et Tempus paschale et festa
  Beatae Virginis et Adventum Domini et festa Christi.
  Hymnus vero \incipit{Salvator mundi}
  diebus sabbatinis dicitur, quando feria die dominico canitur,
  exceptis solemnitatibus.}
}

\rubrica{Sequitur hymnus Sabbatis diebus,
  ad completorium quando ferialiter dicitur.}

\hymnus{salvator_mundi}

\headingCantus

\cantus{hymni_telucis_karlstejn}

\cantus{commune_versiculus}

\cantus{commune_salva1}

\rubrica{Alia nota super eandem antiphonam,
  quae canitur in magnis festis per circulum anni.}\footcite[81r]{xv_a_10}

\cantus{commune_salva2}
